% Invented deck: sections between frames under beamer's default theme, which
% declares no section page, so each one names navigation alone: its anchor and
% label land on the frame set next and its title reaches the footline. A
% subsection, a starred section, references to them, and text set outside every
% frame, which stands on a page of its own as beamer sets it. Nothing here
% derives from any private document.
\documentclass[aspectratio=169]{beamer}

\usetheme{default}
\chrome{ footer = { left = \sectiontitle, right = \framenumber } }

\fonts{
  dir = "fonts",
  body = "Open Sans",
}

\title{Sections Without Pages}
\author{Pat Placeholder}
\date{}

\begin{document}

\maketitle

\section{Kettles}\label{sec:kettles}

\begin{frame}{Filling}
  A kettle fills to its line, and Section~\ref{sec:cups} pours it.
\end{frame}

\subsection{Spouts}\label{sec:spouts}

\begin{frame}{Pouring}
  The spout of subsection~\ref{sec:spouts} sets the pour.
\end{frame}

\section{Cups}\label{sec:cups}

\begin{frame}{Serving}
  Cups follow kettles, as Section~\ref{sec:kettles} said.
\end{frame}

A note set outside every frame stands on a page of its own.

\section*{Washing up}

\begin{frame}{Drying}
  The last frame dries the cups.
\end{frame}

\end{document}
